\documentclass[11pt]{article}

\usepackage[a4paper,margin=1in]{geometry}
\usepackage{amsmath,amssymb,amsthm}
\usepackage{booktabs,enumitem,tabularx}
\usepackage[T1]{fontenc}
\usepackage{microtype}
\usepackage[hypertexnames=false]{hyperref}

\hypersetup{colorlinks=true,linkcolor=blue,citecolor=blue,urlcolor=blue,
  pdftitle={The Sharp Benchmark Constant for the Deng--Hani--Ma Two-Layer Cutting Algorithm},
  pdfauthor={Zhangjinba},pdfsubject={Combinatorial molecule reduction}}

\newtheorem{theorem}{Theorem}[section]
\newtheorem{lemma}[theorem]{Lemma}
\newtheorem{proposition}[theorem]{Proposition}
\newtheorem{corollary}[theorem]{Corollary}

\newcommand{\good}{\#\{33\}}
\newcommand{\bad}{\#\{4\}}
\newcommand{\MU}{M_U}
\newcommand{\MD}{M_D}
\newcommand{\CUTU}{\operatorname{CUT_U}}
\newcommand{\ToyI}{\mathrm{ToyI}}

\title{The Sharp Benchmark Constant for the Deng--Hani--Ma Two-Layer Cutting Algorithm}
\author{Zhangjinba\\Independent Researcher}
\date{August 5, 2026}

\begin{document}
\maketitle

\begin{abstract}\sloppy
The simplified two-layer molecule argument of Deng--Hani--Ma gives
\(\good\ge (|\MU|-1)/5\) for its Toy I cutting algorithm.  We isolate the
source of this loss and prove the stronger, sharp benchmark bound
\[
  \good\ge \left\lceil\frac{|\MU|-1}{3}\right\rceil.
\]
The proof is an amortised flow argument.  Cleanup cuts can orphan cross-bonds,
and operation (b) can orphan one additional cross-bond.  While an upper-layer
atom remains, local slot capacity is paid for by bottom-fixed-end tokens created
by earlier cuts.  An injective token map gives \(\#\mathrm{CUT_U}\le2\good\),
which combines with \(\good=|\MU|-1-\#\mathrm{CUT_U}\).  Three explicit families
give the matching upper bound, so \(a_n=\lceil(n-1)/3\rceil\) for \(n\ge4\).

As a scoped intrinsic consequence, every finite Toy I molecule satisfies
\(v_d(M)\ge\lceil\rho(M)/3\rceil-10d\), and the corresponding Toy I
finite-size constants have asymptotic lower limit at least \(1/3\).  We make no
claim for global \(c_d^*\), larger toy classes, or arbitrary legal sequences.
\end{abstract}

\paragraph{Keywords.} combinatorial molecule reduction; sharp constant;
amortised counting; two-layer graphs; intrinsic cutting value.

\section{Source contract and scope}

We use the two-layer model and Toy I in Section 11 of the public arXiv version
of Deng--Hani--Ma \cite[Definitions 11.1--11.4]{DHM}.  Each atom has four
slots, partitioned into bonds, free ends, and fixed ends; fixed ends do not
count toward degree.  There are at most two parent slots and at most two child
slots.  A cross-layer bond is directed from \(\MU\) to \(\MD\) and occupies a
parent slot of its lower endpoint.  Cutting a connected set as free turns each
boundary bond into one persistent fixed end at the surviving endpoint (bottom
at a surviving parent and top at a surviving child); fixed ends are neither
rewired nor duplicated.  Both layers are trees with no initial fixed ends; no
atom of \(\MD\) is a parent of an atom of \(\MU\); and every atom of \(\MU\)
has exactly one bond to \(\MD\).

The residual \(\MD\) is proper in the source sense.  We run the six-step
algorithm of Definition 11.4, allowing all its tie-breaks.  Proposition 11.7
of the same source gives a complete legal run, elementary output, and
\(\bad=1\).

Call a complete execution of the six-step algorithm, with any permitted
tie-break and an empty final residual molecule, a \emph{Definition-11.4
execution}.  Every cut in such an execution is an elementary component; the
prescribed three-atom \(\{343\}\) cleanup macro is shorthand for its split into
one \(\{3\}\) and one \(\{33\}\), so it contributes one to \(\good\) and none
to \(\bad\).  Separately, a \emph{legal complete cut} means any sequence of
elementary cuts that ends with an empty residual molecule.  The intrinsic value
\(v_d\) maximises the resulting score over these general legal complete cuts,
whereas the benchmark theorem below concerns only Definition-11.4 executions.

For a benchmark run, let \(T=\good\) and \(n=|\MU|\), and define
\[
 a_n=\min_{|\MU|=n}\;\min_{\text{Definition-11.4 tie-breaks}} T,
\]
where the first minimum ranges over finite Toy I molecules.  The source gives
\(\bad=1\) and \(T\ge(n-1)/5\).  Our result replaces only the second estimate.

The intrinsic quantities are
\[
 v_d(M)=\max_C\bigl(\good(C)-10d\,\bad(C)\bigr),\qquad
 W(M)=\max_C\good(C),
\]
where \(C\) ranges over legal complete cuts.  The benchmark theorem addresses
\(a_n\); the scoped transfer to \(v_d\) is in Section~\ref{sec:intrinsic}.

\subsection{Source and related work}

Deng--Hani--Ma supply the original \(1/5\) benchmark estimate and its role in
the hard-sphere derivation \cite{DHM}.  The survey of
Bodineau--Gallagher--Saint-Raymond--Simonella reproduces the surrounding
reduction \cite{BGSS}.  Bruned--Clarisse study a Kruskal-style reduction for a
wave-molecule problem \cite{BC}; that is a structural result on the wave side,
not the quantitative hard-sphere count proved here.  Our claim is therefore a
quantitative Toy I benchmark result, not an algorithm-classification claim or
a new kinetic-limit theorem.  The literature search was bounded, so no
exhaustive priority claim is made.

\section{Exact accounting}

Call operation (a) an \(A\) event, operation (b) a \(B\) event, a cleanup of an
adjacent degree-three pair a \(C\) event, a cleanup of a \(\{343\}\) triple a
\(D\) event, and a one-atom upper-layer removal after its cross-bond partner
has gone a \(U\) event.  Let \(E_<\) be the \(B,C,D\) events whose pre-cut state
still contains at least one upper-layer atom.  Cleanups after the upper layer
is empty are excluded: they cannot orphan a future upper-layer partner.

For a complete run, the counting identities are
\begin{equation}
 T=B+C+D,\qquad A-(C+D)=1,\qquad A+B+U=n. \tag{2.0}
\end{equation}
The first counts the good components produced by (b), \(C\), and the \(D\)
split; the second is the change of the source potential; the third partitions
the upper-layer atoms.

For \(e\in E_<\), let \(S_e\subset\MD\) be its lower-layer cut set.  Define
\(q_e\) as the live cross-bonds removed by \(e\), except that for a \(B\) event
the selected upper--lower bond is omitted.  Define \(s_e\) as the number of
live lower-layer boundary arcs \(p\to v\) with \(p\notin S_e\) and
\(v\in S_e\), counted by edges, and let \(r_e\) be the number of degree-three
atoms in \(S_e\) carrying a bottom fixed end.

Thus \(q_e\) counts future upper atoms orphaned by the event, while each
\(s_e\) creates one bottom-fixed-end token at the surviving lower parent.

\begin{lemma}[Exact orphan partition]
For every complete run,
\begin{equation}
  \#\mathrm{CUT_U}=\sum_{e\in E_<}q_e. \tag{2.1}
\end{equation}
\end{lemma}

\begin{proof}
Every \(\mathrm{CUT_U}\) atom has lost its unique initial cross-bond partner
earlier.  A \(C\) or \(D\) event removes exactly the live cross-bonds counted
by \(q_e\).  A \(B\) event removes its selected upper atom together with the
selected bond; its other live cross-bonds are exactly those counted by \(q_e\).
Conversely, every counted bond has a still-present upper endpoint and produces
one later \(\mathrm{CUT_U}\).  Events after the upper layer is empty have no
live cross-bond to a future upper atom.\end{proof}

\begin{lemma}[Local flow capacity]
For every \(e\in E_<\),
\begin{equation}
 q_e+s_e\le2+r_e. \tag{2.2}
\end{equation}
\end{lemma}

\begin{proof}
For a degree-three lower atom \(v\), let \(c_v\) be its live cross-bonds,
\(p_v\) its live lower-layer parents, and let \(t_v\in\{0,1\}\) indicate
whether its unique fixed end is bottom.  The atom has two parent slots; a top
fixed end occupies one of them and a bottom fixed end does not.  Hence
\begin{equation}
 c_v+p_v+(1-t_v)\le2. \tag{2.3}
\end{equation}
For \(B\), \(q_e=c_v-1\), \(s_e=p_v\), and \(r_e=t_v\), giving the stronger
\(q_e+s_e\le r_e\).  For \(C\), one internal lower bond contributes one
parent slot, so summing (2.3) gives \(q_e+s_e\le1+r_e\).  For \(D\), the two
outer atoms have degree three, the centre has degree four and no fixed end, and
the two internal bonds contribute two parent slots.  Summing gives
\(q_e+s_e\le2+r_e\).\end{proof}

\begin{lemma}[Bottom-token injection]
\begin{equation}
 \sum_{e\in E_<}r_e\le\sum_{e\in E_<}s_e. \tag{2.4}
\end{equation}
\end{lemma}

\begin{proof}
Initially the lower layer has no fixed ends.  Before an event in \(E_<\), the
only operations removing lower atoms are earlier \(B,C,D\) events.  Operations
(a) and \(U\) remove only upper atoms and create top, not bottom, fixed ends on
the lower side; step (6) occurs only after the upper layer is empty.  Thus each
bottom fixed end counted by an \(r_e\) has a unique earlier boundary arc counted
by some \(s_f\).  The fixed end is never rewired or duplicated, so distinct
consumed tokens have distinct source arcs.\end{proof}

\begin{theorem}[Sharp Toy I benchmark bound]
For every finite Toy I molecule and every legal tie-break of Definition 11.4,
\begin{equation}
 \good\ge\left\lceil\frac{|\MU|-1}{3}\right\rceil. \tag{2.5}
\end{equation}
\end{theorem}

\begin{proof}
Summing (2.2) and using (2.4) gives
\[
 \sum q_e\le2|E_<|\le2\good.
\]
The first inequality uses cancellation of the \(s_e\) terms, and the second
uses that every \(B,C,D\) event in \(E_<\) produces one good component.  By
(2.1), \(\#\mathrm{CUT_U}\le2\good\).  From (2.0), eliminating the event
counts gives
\[
 \good=|\MU|-1-\#\mathrm{CUT_U}.
\]
Therefore \(|\MU|-1\le3\good\), and integrality gives (2.5).\end{proof}

The exact run-level contract is
\[
 \sum_{\text{cleanup }k}b_k+\sum_{B\text{ events }e}q_e
 \le2(B+C+D).
\]
The older condition \(\sum b_k\le2(C+D)\) is a stronger sufficient condition;
it is neither used nor claimed here.

\section{Sharpness: matching explicit families}

The three families below are finite Toy I molecules.  Their lower layers are
paths and their upper layers are paths; the cross-bond maps are given explicitly
so that the construction is independent of enumeration.

\paragraph{Transpose family.}
For \(k\ge1\), take \(n=3k+1\), let \(\MU=(u_0,\ldots,u_{3k})\) and
\(\MD=(d_0,\ldots,d_{3k-1})\) be directed paths with internal bonds
\(u_i\to u_{i+1}\) and \(d_j\to d_{j+1}\), and for
\(0\le i<3k\) write \(i=qk+r\), \(q\in\{0,1,2\}\), \(0\le r<k\), and join
\(u_i\) to \(d_{3r+q}\).  Add \(u_{3k}\to d_0\).  The lowest upper atoms
are removed in the order \(u_{3k},u_{3k-1},\ldots,u_{2k}\), giving \(k+1\)
operations (a).  Their lower partners are \(d_0,d_{3r+2}\).  After these
cuts the degree-three set is \(\{d_0\}\cup\{d_{3r+2}:0\le r<k\}\); the
points are separated along the path until the final cut, when \(d_0,d_2\)
and their degree-four neighbour \(d_1\) form the first \(\{343\}\).  The
blocks \(\{d_{3r},d_{3r+1},d_{3r+2}\}\) can then be removed successively in
the indicated legal tie-break.  For \(0\le r<k-1\), after block \(r\) is
removed, \(d_{3r+3}\) becomes degree three,
\(d_{3r+5}\) was already degree three, and \(d_{3r+4}\) remains degree four
because its cross-bond comes from the still present \(u_{k+r+1}\).  In block
\(r\), the three designated cross-bonds come from
\(u_r,u_{k+r},u_{2k+r}\).  For \(r\ge1\), only the last of these endpoints
has already been removed, so the cleanup breaks two live cross-bonds.  For
\(r=0\), there is in addition \(u_{3k}\to d_0\), but both this endpoint and
the designated endpoint \(u_{2k}\) have already been removed; again exactly
two live cross-bonds are broken.  The final block \(r=k-1\) is the terminal
cleanup, so no successor indices are needed.  This proves inductively that the
block cascade is the stated one.  After the \(k\) cleanups, the remaining \(2k\)
upper atoms are removed by \(\CUTU\).  Thus
\[
 (A,B,C,D,U)=(k+1,0,0,k,2k),\qquad \good=k.
\]

\paragraph{Padded family.}
Starting from the transpose family, append one atom to the top of \(\MU\) and
one atom to the end of \(\MD\), and join the two new atoms.  The old lower
endpoint changes from two bonds and two free ends to three bonds and one free
end, so its degree is unchanged.  The new upper endpoint is removed last;
therefore the transpose phase and all its \(k\) cleanups are unchanged.  The
last lower atom then has degree three, and its pair with the remaining upper
atom is one operation (b).  For \(n=3k+2\),
\[
 (A,B,C,D,U)=(k+1,1,0,k,2k),\qquad \good=k+1.
\]

\paragraph{Trimmed family.}
For \(k\ge2\), take \(\MU=(u_0,\ldots,u_{3k-1})\) and
\(\MD=(d_0,\ldots,d_{3k-2})\) as paths.  Use the transpose cross-bonds except
that the bond which would end at \(d_{3k-1}\) is redirected to \(d_0\); the
extra \(u_{3k}\) is omitted.  The first \(k\) upper cuts and the next
\(k-1\) block cleanups are the same as above.  The cascade stops at the tail
pair \(d_{3k-3},d_{3k-2}\): one further operation (a) lowers the latter to
degree three, after which the pair is a \(\{33\}\) cleanup.  This final cleanup
breaks one live cross-bond, while each preceding block breaks two.  The
remaining \(2k-1\) upper atoms are then \(\CUTU\) atoms, and
\[
 (A,B,C,D,U)=(k+1,0,1,k-1,2k-1),\qquad \good=k.
\]

Together with Theorem 2.1, these families prove
\begin{equation}
 a_n=\left\lceil\frac{n-1}{3}\right\rceil,\qquad n\ge4. \tag{3.1}
\end{equation}

\section{Intrinsic consequence on Toy I}\label{sec:intrinsic}

The benchmark theorem supplies one legal witness in the intrinsic maximum; it
does not assert that every legal cutting sequence has the same count.

\begin{proposition}[Scoped intrinsic lower bound]
Fix \(d>0\).  For every finite Toy I molecule with both layers nonempty,
\begin{equation}
 v_d(M)\ge\left\lceil\frac{|\MU|-1}{3}\right\rceil-10d
 =\left\lceil\frac{\rho(M)}3\right\rceil-10d. \tag{4.1}
\end{equation}
\end{proposition}

\begin{proof}
Let \(n=|\MU|\) and \(m=|\MD|\).  The two layer trees have \(n-1\) and
\(m-1\) internal edges, and Toy I has exactly \(n\) cross-bonds.  The full
graph is connected.  Here \(\rho(M)\) denotes its circuit rank,
\(\rho(M)=|E|-|V|+1\), so
\[
 \rho(M)=|E|-|V|+1=((n-1)+(m-1)+n)-(n+m)+1=n-1.
\]
The source algorithm supplies a legal complete run with \(\bad=1\), and
Theorem 2.1 gives \(\good\ge\lceil(n-1)/3\rceil\) for it.  Since \(v_d\) is
the maximum score over all legal complete cuts, it is at least the score of
this witness.\end{proof}

For \(n\ge4\), define the finite-size constant by an infimum (the Toy I
class has no a priori bound on \(|\MD|\)):
\[
 c_d^{\ToyI}(n)=\inf_{\substack{M\text{ finite Toy I}\\|\MU|=n}}
 \frac{v_d(M)}{\rho(M)}.
\]
Then
\[
 c_d^{\ToyI}(n)\ge
 \frac{\lceil(n-1)/3\rceil-10d}{n-1},
 \qquad
 \liminf_{n\to\infty}c_d^{\ToyI}(n)\ge\frac13. \tag{4.2}
\]
This is an asymptotic statement within Toy I, not the global infimum \(c_d^*\).

\section{Audit and reproducibility}

The proof uses no finite enumeration.  A bounded full-labelled replay checks
11,590 molecules and 25,243 complete tie-break records with total size at most
six, including double-parent configurations in both layers.  It reports no
failed records and no violations of the local flow contract.  This replay is
regression evidence only.

The public code and reproducibility repository is available at:
\begin{center}
\url{https://github.com/xingzhicn/molecule-cut}
\end{center}
It reproduces the exact subset-DP values and the family profiles.  The
reproducibility commands are
\begin{verbatim}
uv run --frozen pytest
uv run --frozen ruff check .
\end{verbatim}
On August 5, 2026 these checks reported 652 passing tests, a clean Ruff check,
and all source hashes valid.

\section{What is and is not concluded}

We prove the exact Toy I/Definition-11.4 benchmark constant \(1/3\), its explicit
sharpness families, and the scoped intrinsic consequences (4.1)--(4.2).  We do
not prove a global positive lower bound for \(c_d^*\), a theorem for \(W\), a
result for Toy I plus/II/III or the full molecule, a score bound for every
arbitrary legal sequence, or any complexity-hardness classification.  No claim
is made about global no-localization, no-beam, or impossibility of the broader
DHM/Boltzmann program.

\begin{thebibliography}{9}

\bibitem{DHM}
Y.~Deng, Z.~Hani, and X.~Ma,
\emph{Long time derivation of the Boltzmann equation from hard sphere dynamics},
arXiv:2408.07818v3 (2025), especially Section 11.2 and Proposition 11.7.
\url{https://arxiv.org/abs/2408.07818}.

\bibitem{BGSS}
T.~Bodineau, I.~Gallagher, L.~Saint-Raymond, and S.~Simonella,
\emph{Derivation of the Boltzmann equation from hard-sphere dynamics
(after Deng, Hani, Ma)}, S\'eminaire Bourbaki 1247, arXiv:2602.04407, 2026.

\bibitem{BC}
Y.~Bruned and V.~Clarisse,
\emph{Kruskal-style algorithm for cubic Schr\"odinger equation molecule
reduction}, arXiv:2603.23298, 2026.

\end{thebibliography}

\end{document}
